\section{Track-to-vertex association with graph networks}
\label{sec:gnn}

\subsection{Motivation and task}
\label{sec:gnn:task}

Vertex finding answers \emph{where} the primary vertices are but not \emph{which
tracks belong to which vertex}. Physics analyses need the second answer: the
hard-scatter vertex is chosen as the vertex with the largest
$\sum p_{\mathrm{T}}^{2}$ of its associated tracks, pile-up jet rejection cuts on
track-to-vertex association, and any per-vertex quantity (multiplicity, mass,
displaced-object pointing) requires a track list. AMVF produces these lists as a
by-product of its fit; \pvf{} does not, so a dedicated Track-to-Vertex
Association (TTVA) step completes the chain,
\begin{center}
tracks $\;\to\;$ \pvf{} histogram $\;\to\;$ peak finding $\;\to\;$
candidate vertices $\;\to\;$ \textbf{TTVA} $\;\to\;$ per-vertex track lists.
\end{center}
Assigning tracks also lets \pvf{} vertices be judged by the same track-based
purity criteria applied to AMVF, giving an apples-to-apples comparison of the two
chains end to end (Section~\ref{sec:gnn:metrics}).

TTVA is formulated as \emph{binary edge classification} on a per-event graph: a
score for every (track, vertex) pair, one selection rule on top. A graph neural
network is a natural fit because the decision for one track depends on the other
tracks and vertices in its neighbourhood (a vertex that already attracts many
compatible tracks should attract this one more strongly), which fixed per-pair
classifiers cannot express. The key design decision, following the November 2025
prototype, is to \textbf{train on Monte-Carlo truth vertices} (where edge labels
exist) and \textbf{deploy on \pvf{} peaks}. The gap this opens between training
and deployment conditions, and how it was closed, is the subject of
Section~\ref{sec:gnn:transfer}.

\subsection{Graphs, features and labels}
\label{sec:gnn:graphs}

Each event is a heterogeneous bipartite graph (Figure~\ref{fig:gnn_graph}).
\emph{Track nodes} carry eight features: \dzero{}, \zzero{}, their
uncertainties and covariance, the polar and azimuthal angles, and
$p_{\mathrm{T}}/50\,\mathrm{GeV}$. \emph{Vertex nodes} carry two: the $z$
position and the histogram peak height (for truth graphs, a height synthesized
by the same Gaussian-CDF recipe that builds the \pvf{} targets of
Section~\ref{sec:method:targets}). Each track--vertex edge carries three
physics attributes: the longitudinal significance
$(\zzero-z_{\mathrm{PV}})/\sqrt{\sigma_{z_0}^{2}+\sigma_{\mathrm{PV}}^{2}}$, the
transverse significance $\dzero/\sigma_{d_0}$, and $|\zzero-z_{\mathrm{PV}}|$.
Training labels mark an edge as true iff MC truth associates the track with that
vertex.

At \muavg{}$\approx$60 the graphs are fully connected (every track paired with
every vertex). At PU200 a fully connected event has $\sim$118\,000 edges
($\sim$930 tracks $\times$ 128 vertices), so each track is connected only to its
$k=20$ nearest vertices in $|\Delta z|$; this retains $99.5\%$ of true edges
(measured on 200 events; the lost $0.5\%$ are tracks whose \zzero{} lies tens of
millimetres from their true vertex, unlearnable for any $z$-based method) and
keeps the edge count at $\sim$18\,500.

\begin{figure}[t]
  \centering
  \resizebox{0.9\linewidth}{!}{%
  \begin{tikzpicture}[x=1cm,y=1cm,>=Latex,font=\small]
    % track nodes
    \foreach \y/\n in {4.4/1, 3.3/2, 2.2/3, 1.1/4} {
      \filldraw[fill=mydarkblue!14, draw=mydarkblue!70, thick] (0,\y) circle (10pt)
        node[font=\scriptsize, text=mydarkblue!90!black] {$t_\n$};
    }
    \node[align=center, font=\scriptsize, text=black!60] at (0,5.4)
      {track nodes\\$(\dzero,\zzero,\sigma_{d_0},\sigma_{z_0},\mathrm{cov},\theta,\phi,p_{\mathrm{T}})$};
    % pv nodes
    \foreach \y/\n in {3.85/1, 2.2/2, 0.55/3} {
      \filldraw[fill=myorange!18, draw=myorange!80!black, thick] (6,\y) circle (10pt)
        node[font=\scriptsize, text=myorange!60!black] {$v_\n$};
    }
    \node[align=center, font=\scriptsize, text=black!60] at (6,5.4)
      {vertex nodes\\$(z,\ \text{peak height})$};
    % edges
    \foreach \s/\t in {4.4/3.85, 4.4/2.2, 3.3/3.85, 3.3/2.2, 2.2/3.85, 2.2/2.2, 2.2/0.55, 1.1/2.2, 1.1/0.55} {
      \draw[black!45, thick] (0.37,\s) -- (5.63,\t);
    }
    \draw[accent, very thick] (0.37,3.3) -- (5.63,2.2);
    \node[align=center, font=\scriptsize, text=accent] at (3.0,3.6)
      {edge attributes:\\
       $\dfrac{\zzero-z_{\mathrm{PV}}}{\sqrt{\sigma_{z_0}^{2}+\sigma_{\mathrm{PV}}^{2}}},\quad
        \dfrac{\dzero}{\sigma_{d_0}},\quad |\zzero-z_{\mathrm{PV}}|$};
    % model
    \node[draw=mydarkblue!60, rounded corners=8pt, minimum width=2.6cm, minimum height=1.15cm,
          align=center, fill=mydarkblue!12, line width=0.9pt] (gat) at (9.6,3.4)
      {$2\times$ hetero GAT\\ \scriptsize (4 heads, edge-aware)};
    \node[draw=mydarkblue!60, rounded corners=8pt, minimum width=2.6cm, minimum height=1.15cm,
          align=center, fill=mydarkblue!12, line width=0.9pt] (mlp) at (9.6,1.4)
      {edge MLP\\ \scriptsize $[64{\to}32{\to}32{\to}1]$};
    \node[draw=myorange!70!black, rounded corners=6pt, minimum width=2.6cm, minimum height=1.0cm,
          align=center, fill=myorange!16, line width=0.9pt] (score) at (9.6,-0.4)
      {edge score\\ \scriptsize $\in[0,1]$ per (track, vertex)};
    \draw[->, thick, draw=black!60] (7.2,3.4) -- (gat);
    \draw[->, thick, draw=black!60] (gat) -- node[right, font=\tiny, text=black!55]{node embeddings} (mlp);
    \draw[->, thick, draw=black!60] (mlp) -- (score);
  \end{tikzpicture}}
  \caption{TTVA graph and model. Tracks and candidate vertices are the two node
  types of a bipartite graph whose edges carry impact-parameter significances
  and $|\Delta z|$; two heterogeneous graph-attention layers refine the node
  embeddings, and a small MLP scores each edge. At PU200 each track connects to
  its 20 nearest vertices in $|\Delta z|$.}
  \label{fig:gnn_graph}
\end{figure}

\subsection{Architecture and training}
\label{sec:gnn:arch}

Track and vertex features are first embedded into 32 dimensions by separate
linear layers. Two rounds of heterogeneous message passing follow, each a pair
of graph-attention convolutions~\cite{velivckovic2017graph} (one per edge
direction, four heads, max aggregation, residual connections) whose attention
consumes the three edge attributes, so the significances bias the messages
directly. A three-layer MLP on the concatenated (track, vertex) embeddings
yields one logit per edge. The model has $39\,521$ parameters --- $90\times$
smaller than the \pvf{} finder itself --- and is trained with
binary cross-entropy whose positive-class weight is recomputed per batch
(most pairs are not associated), Adam at $10^{-3}$, batches of 32 event graphs.
The final training recipe adds cosine learning-rate annealing to $10^{-5}$ and
gradient-norm clipping at $1.0$: at fixed learning rate the PU200 training loss
collapsed and recovered between epochs 50 and 75, while the annealed run is
smooth throughout and ends $9.5\%$ lower in validation loss.

Three trainings enter this note: the November 2025 baseline (51k
\muavg{}$\approx$60 $t\bar{t}$ events; reproduced bit-exactly by the current
package), a 30k-event PU200 training on the single-lepton ITk sample, and the
final PU200 training on 180\,000 all-hadronic events
(Section~\ref{sec:gnn:transfer}) with chain-like augmentation, drawn from the
$2.5\times10^{6}$-event sample of Section~\ref{sec:data:hllhc} and cycled
through memory in 20k-graph shards.

\subsection{Selection, metrics and vertex conventions}
\label{sec:gnn:metrics}

At inference each track is assigned to at most one vertex: the vertex with the
highest edge score, kept only if that score exceeds a threshold $t$
(\emph{MaxScore} selection); otherwise the track stays unassigned. The model is
run once and $t$ is a free post-hoc knob. Scores saturate near 1 for
unambiguous tracks, so useful operating points live at $t\geq0.9$ and the scan
must extend to $t=0.999$: at PU200 moving from the naive $t=0.5$ to the working
point is worth $+11$ points of the headline metric with no retraining.

Reconstructed vertices are then classified by the truth composition of their
assigned tracks: \emph{Clean} if $\geq50\%$ of assigned tracks come from one
truth vertex, \emph{Merged} otherwise, \emph{Split} if another reconstructed
vertex already claimed the same dominant truth vertex with larger
$\sum p_{\mathrm{T}}^{2}$, and \emph{Fake} if no assigned track has a truth
vertex. The headline metric is the \textbf{clean-vertex yield}
$N_{\text{clean}}/N_{\text{truth}}$ with $N_{\text{truth}}$ the truth vertices
with $\ntrk\geq2$. Two remarks keep this comparable and honest. First, this
\emph{track-based} taxonomy is stricter than, and not numerically comparable
to, the position-based matching of Section~\ref{sec:method:metrics} used for
the finder-only results; AMVF is always evaluated through the identical
classifier, reading its track lists from the same ntuples. Second, a vertex
with \emph{no} assigned tracks is not a reconstructed vertex in any deployable
sense --- no fit could be performed --- so trackless candidates are dropped
before classification. The historical convention instead booked each one as
Fake; under that bookkeeping a model that correctly abstains on a junk peak is
punished exactly like one that hallucinates a vertex. The drop-empty convention
leaves the clean-vertex yield untouched (an empty vertex is never Clean) and
moves only the fake accounting; both conventions are reported in the
repository outputs.

\subsection{Low-pileup performance and real-data validation}
\label{sec:gnn:mu60}

On the 2\,550-event \muavg{}$\approx$60 test sample the edge classifier is
essentially solved: ROC AUC $0.998$ over $45.3$ million labelled edges, with
the score distributions of true and false edges separated by three orders of
magnitude at the extremes (Figure~\ref{fig:gnn_mu60}, top left). The
threshold scan (top right) sets the working point $t^{*}=0.98$, where the
clean-vertex yield of the full \pvf{}+GNN chain reaches $\mathbf{74.5\%}$ at a
$0.9\%$ fake rate, against $\mathbf{60.9\%}$ for AMVF judged identically
(Table~\ref{tab:gnn_mu60}). The efficiency-vs-multiplicity curve (bottom left)
shows the chain above AMVF in every \ntrk{} bin, with the largest gains for
low-multiplicity vertices, which AMVF tends to absorb into merged fits (its
merged rate is $20\%$ against $11.8\%$ for the chain). As a pure associator the
GNN also edges AMVF at track level: efficiency$\times$purity $F_1 = 0.866$ vs
$0.849$. Hard-scatter identification is statistically tied ($98.1\%$ vs
$98.9\%$).

\begin{table}[h]
  \centering
  \caption{Track-based vertex classification at \muavg{}$\approx$60
  (2\,550 test events, 72\,189 truth PVs with $\ntrk\geq2$), MaxScore
  $t^{*}=0.98$. All rows judged by the identical classifier. The last row uses
  ground-truth vertex positions as graph nodes and bounds the associator alone.}
  \label{tab:gnn_mu60}
  \begin{tabular}{lcccccc}
    \toprule
    Chain & Reco PVs & Clean & Merged & Split & Fake & Clean/truth \\
    \midrule
    \pvf{} peaks + GNN & 66\,472 & $80.9\%$ & $11.8\%$ & $6.3\%$ & $0.9\%$ & $\mathbf{74.5\%}$ \\
    AMVF               & 57\,329 & $76.7\%$ & $20.0\%$ & $2.5\%$ & $0.7\%$ & $60.9\%$ \\
    truth vertices + GNN (bound) & 80\,498 & $79.9\%$ & $7.6\%$ & $3.9\%$ & $8.6\%$ & $89.1\%$ \\
    \bottomrule
  \end{tabular}
\end{table}

\begin{figure}[p]
  \centering
  \begin{subfigure}{0.49\linewidth}
    \includegraphics[width=\linewidth]{gnn_edge_scores.png}
    \caption{Edge scores by truth label.}
  \end{subfigure}\hfill
  \begin{subfigure}{0.49\linewidth}
    \includegraphics[width=\linewidth]{gnn_mu60_scan.png}
    \caption{Working-point scan.}
  \end{subfigure}
  \par\medskip
  \begin{subfigure}{0.49\linewidth}
    \includegraphics[width=\linewidth]{gnn_mu60_eff_ntrk.png}
    \caption{Clean-vertex efficiency vs \ntrk.}
  \end{subfigure}\hfill
  \begin{subfigure}{0.49\linewidth}
    \includegraphics[width=\linewidth]{gnn_mu60_categories.png}
    \caption{Category rates at $t^{*}$.}
  \end{subfigure}
  \caption{TTVA at \muavg{}$\approx$60. The edge classifier is near-saturated
  (a); the clean-vertex yield rises monotonically up to $t^{*}=0.98$ before
  fake vertices take over (b); the \pvf{}+GNN chain beats AMVF in every truth
  multiplicity bin (c), and trades AMVF's merges for a higher clean fraction
  (d).}
  \label{fig:gnn_mu60}
\end{figure}

The same chain runs unchanged on Run~3 collision data, where no truth exists
and AMVF is used as a reference instead: $84.6\%$ of AMVF-assigned tracks land
on the GNN vertex matched to their AMVF vertex ($92.6\%$ among tracks both
methods assign), the leading-$\sum p_{\mathrm{T}}^{2}$ vertex agrees in
$84.8\%$ of events, and an AMVF-vs-AMVF closure through the same machinery
gives exactly $1.0$. The edge-score distributions on data closely track the
simulation, with a $\sim3\times$ larger population of ambiguous tracks at
intermediate scores --- a mild, visible domain shift consistent with the larger
unassigned fraction on data.

\subsection{HL-LHC: the full chain against AMVF}
\label{sec:gnn:pu200}

The PU200 evaluation runs the production v4b finder (Section~\ref{sec:hllhc})
at its production operating point (amplitude $10^{-2}$, integral $0.40$,
minimum width 3, height floor $0.03$) on 1\,500 held-out single-lepton events,
builds $k=20$ graphs on the found peaks ($88.4$ per event), and scores them
with the GNN; AMVF track lists from the same ntuple go through the identical
classifier (148\,133 truth PVs with $\ntrk\geq2$).

Decomposing the achievable yield locates every loss
(Table~\ref{tab:gnn_pu200}, Figure~\ref{fig:gnn_pu200}). The finder places a
peak within $0.5$~mm of $81.0\%$ of truth vertices; the missed ones are
low-multiplicity (median $\ntrk=3$ against $7$ for found vertices).
Re-assigning every track to the peak nearest its \emph{true} vertex --- a
perfect-information oracle --- yields $74.8\%$: the $6$-point drop from the
finder cap is close vertex pairs sharing a single nearest peak, unrecoverable
by any associator on these peaks ($11.9\%$ of oracle vertices are Merged).
The GNN reaches $64.7\%$ of truth at its working point $t=0.995$ with a fake
rate of $0.08\%$ --- $86\%$ of the oracle --- while AMVF, whose own finder and
fitter face the same pile-up density, delivers $57.3\%$ at an $0.91\%$ fake
rate. The \pvf{}+GNN chain therefore outperforms AMVF by $7.4$ points of
clean-vertex yield at an eleven-fold lower fake rate. Hard-scatter
identification through the chain ties AMVF at $97.5\%$ (both at $t=0.5$;
HS-ID and vertex classification use different thresholds on the same forward
pass, $92.1\%$ if the clean-vertex threshold is reused). %% V3: add final v3 row + sentence after evaluation

\begin{table}[h]
  \centering
  \caption{PU200 full chain against AMVF (1\,500 events, 148\,133 truth PVs,
  drop-empty convention). ``Finder cap'' counts truth PVs with a peak within
  $0.5$~mm (greedy one-to-one); the oracle assigns every track to the peak
  nearest its true vertex and bounds any associator on these peaks. GNN rows
  use the 30k-event PU200 training (v1).}
  \label{tab:gnn_pu200}
  \begin{tabular}{llcc}
    \toprule
    & Configuration & Clean/truth & Fake rate \\
    \midrule
    reference & AMVF & $57.3\%$ & $0.91\%$ \\
    chain & \pvf{} v4b + GNN, $t=0.95$ & $58.0\%$ & $0.08\%$ \\
    chain & \pvf{} v4b + GNN, $t=0.995$ & $\mathbf{64.7\%}$ & $0.08\%$ \\
    %% V3 row here after evaluation
    bound & oracle association on v4b peaks & $74.8\%$ & $0.00\%$ \\
    bound & finder cap ($0.5$~mm) & $81.0\%$ & --- \\
    bound & truth vertices + GNN & $91.8\%$ & --- \\
    \bottomrule
  \end{tabular}
\end{table}

\begin{figure}[p]
  \centering
  \begin{subfigure}{0.9\linewidth}
    \centering
    \includegraphics[width=\linewidth]{pu200_yield_ladder.png}
    \caption{Clean-vertex yield with the three bounds.}
  \end{subfigure}
  \par\medskip
  \begin{subfigure}{0.55\linewidth}
    \includegraphics[width=\linewidth]{pu200_chain_scan.png}
    \caption{Threshold scan (drop-empty).}
  \end{subfigure}\hfill
  \begin{subfigure}{0.43\linewidth}
    \includegraphics[width=\linewidth]{pu200_miss_ntrk.png}
    \caption{Finder miss probability vs \ntrk.}
  \end{subfigure}
  \caption{PU200 full chain. (a) The \pvf{}+GNN chain exceeds AMVF by
  $7.4$ points of clean-vertex yield; the dashed bounds separate what the
  associator can still recover (up to the oracle) from what belongs to the
  finder (oracle to cap) and to pile-up density itself. (b) The yield rises
  monotonically to $t\approx0.995$ while the drop-empty fake rate stays an
  order of magnitude below AMVF's at every threshold. (c) The finder misses
  are concentrated at low track multiplicity.}
  \label{fig:gnn_pu200}
\end{figure}

\subsection{The truth-to-peak transfer gap}
\label{sec:gnn:transfer}

Training on truth vertices and deploying on peaks opens a distribution gap on
the vertex side of the graph, and three concrete mechanisms were identified.
First, a bookkeeping artefact: under the all-peaks convention the
better-converged cosine-annealed model appeared to have a $4.5\%$ chain fake
floor against $0.6\%$ for the earlier fixed-rate model --- but the excess was
entirely junk peaks the model correctly left \emph{empty}, booked as Fake; the
gap vanishes under the drop-empty convention
(Section~\ref{sec:gnn:metrics}). Second, a genuine bug inherited from the
original prototype: the Gaussian-CDF height recipe evaluated the CDF $240$~mm
away from every vertex (a beamline offset applied twice), so \emph{every
truth-training graph ever built carried vertex heights identically zero} while
deployment graphs carry real peak heights up to $\sim4$. The weights
multiplying the height input were unconstrained at training time and injected
uncontrolled contributions at inference --- and better-optimized checkpoints
can be \emph{more} sensitive to the dead input, explaining why the improved
training transferred worse. Third, distribution shift in the remaining vertex
quantities, measured directly on the chain: peak positions differ from truth
$z$ by the finder residuals, peak widths (which enter the edge significances)
are $2$--$5\times$ the preset resolution model, $9.4\%$ of peaks are junk with
a characteristic low-height spectrum, and $19\%$ of truth vertices have no
peak at all --- their tracks must be abstained on, a situation truth graphs
never show.

The final training closes the gap at the source: the height recipe is fixed,
and truth graphs are made \emph{chain-like} at build time with probability
$0.7$ per event --- vertices are dropped with the measured miss probability
versus \ntrk{}, surviving positions are jittered by the measured
peak--truth residuals, preset widths and recipe heights are replaced by
samples of the measured peak-width and height-ratio distributions, and junk
vertex nodes are injected at the measured rate with the measured height
spectrum and all-false edges. Every distribution is an empirical quantile grid
measured on the v4b chain itself, so the augmentation has no free parameters.
Trained on 180k augmented all-hadronic events, this is the v3 model of
Table~\ref{tab:gnn_pu200}. %% V3: replace placeholder sentence with results

\subsection{Computational cost}
\label{sec:gnn:latency}

Per-event wall times at \muavg{}=200 (single A100, batch 1, medians over
1\,500 events) are listed in Table~\ref{tab:gnn_latency}. In the first
implementation the chain spent $74\%$ of its time in the pure-Python peak
scan; JIT-compiling the identical algorithm (verified bit-identical on 300
real histograms) removes it from the budget, and replacing the per-track
selection loop with a vectorized group-maximum (verified to differ only at
exact score ties) removes most of the rest. The optimized chain runs in
$\sim$11~ms per event, of which the GNN forward pass is $3.8$~ms --- the
association step adds a third to the finder's cost and is never the
bottleneck. Both learned components are pile-up-independent in depth; only
the graph size grows with occupancy.

\begin{table}[h]
  \centering
  \caption{Full-chain latency per event at PU200 (medians, single A100,
  batch 1). The optimized peak finder and selection are verified bit-identical
  (peaks) and tie-only-differing (selection) against the originals.}
  \label{tab:gnn_latency}
  \begin{tabular}{lcc}
    \toprule
    Stage & original [ms] & optimized [ms] \\
    \midrule
    \pvf{} forward (12 intervals) & $4.2$ & $4.2$ \\
    peak finding                  & $60.8$ & $0.07$ \\
    graph construction            & $2.0$ & $2.0$ \\
    GNN forward                   & $3.8$ & $3.8$ \\
    MaxScore selection            & $11.7$ & $0.9$ \\
    \midrule
    total                         & $82.5$ & $\sim11$ \\
    \bottomrule
  \end{tabular}
\end{table}

\subsection{Outlook}
\label{sec:gnn:outlook}

The decomposition of Table~\ref{tab:gnn_pu200} prescribes the next steps in
order of leverage. The $19\%$ of truth vertices the finder misses (low
multiplicity) and the $6\%$ lost to peaks shared by close pairs belong to the
finder, reinforcing the conclusions of Section~\ref{sec:hllhc:failure}; an
associator-aware finder objective (emitting separate peaks when the
downstream track assignment could distinguish them) is the natural joint
target. On the associator side the remaining headroom to the oracle is
$\sim$10 points. The all-hadronic PU200 sample offers $2.5\times10^{6}$
training events of which $180$k are used here. HGTD timing remains a
forward-track feature: $4.0\%$ of PU200 tracks carry a valid time overall,
but $44.6\%$ of tracks with $|\eta|>2.3$ do ($8.9\%$ of all tracks), so a
time-augmented edge significance is expected to pay off specifically where
$z$-ambiguities are worst --- a targeted forward study rather than a global
gain (Section~\ref{sec:negative:timing}). Finally, the track-based
classification used throughout this section is the deployment-faithful
metric for the combined system and is the recommended headline for any
future \pvf{}-vs-AMVF comparison: per-task working points ($t=0.5$ for
hard-scatter identification, $t=0.995$ for vertex classification at PU200,
$t^{*}=0.98$ at \muavg{}$\approx$60) come for free from a single forward
pass.
